\pagebreak

\section{Conversion of 3PL struct and enum types to C types - ctypes.3pl}\label{chap:ctypes}

    \index{library!ctypes.3pl}
    This library contains two functions. The first takes a 3PL struct or
    enumerated type and generates an equivalent C type as a string.
    The second function creates a string array of identifiers from an
    enumerated type.
    
    This library must be incorporated by using a {\bf module} statement, i.e.
    \begin{lstlisting}[gobble=4, language=threepl]
       module "ctypes.3pl"
    \end{lstlisting}



    \subsection{library modules}
    
    
        None.


    \subsection{library procedures}
    
        None.




    \subsection{library functions}
    



        \subsubsection{ctypes - converts 3PL types to C types}\label{sec:lfctypes}

            {\bf Library: ctypes.3pl}

            {\bf Prototype:} \\
            \vsp
            \begin{lstlisting}[gobble=12, language=threepl]
               global function
               ctypes (type(t), str(name), str(prefix, "")) {}
            \end{lstlisting}

            {\bf Return mode and type:} \\
            immediate, "str"

            {\bf Input parameters:} \\            
            {\bf t} is the type to be converted. \\
            {\bf name} is the identifier to be used in the C type.\\
            {\bf prefix} is an optional string prefix to prepend to enum
            identifiers.

            {\bf description:} \\
            This function takes a 3PL variable of type {\bf type} and
            returns a string representing the equivalent C type. It is
            intended to allow non-primitive data to be passed between an FPGA and
            a CPU via a 32 bit port. The conversion allows types declared in 3PL
            to be converted to C types for inclusion and compilation.
            
            The only types converted are {\bf struct} and {\bf enum}. A struct may
            contain enums but must not contain nested structs. Any enums contained
            within a struct must have been previously converted by this function.
            Enums where the ordinals have been specified explicitly and
            non-sequentially are converted correctly.
            
            The second argument is the identifier to be given to the enum or
            struct. Since the type is passed by value from the argument there is
            no type identifier available, so one must be provided.
            
            3PL structs are converted to C bit field structs with an additional
            field to pack the struct to 32 bits if necessary. If the total
            number of bits is greater than 32 an error exit occurs. The packed
            bit fields are in reverse order as required.

            As an example, the following 3PL code -
            \begin{lstlisting}[gobble=12, language=threepl]
               type(grptype, "[16]log");
               enum(src_type,
                       TIME,
                       PIXEL,
                       CPU
               );
               enum(subj_type,
                       NONE,
                       LINEARISE,
                       GAINTRIM,
                       GROUP,
                       PILEUP,
                       THROTTLE,
                       ENERGY,
                       DACOLMAP,
                       TIMER,
                       TIMER_IM 
               );
               struct(command,
                   src,        src_type,
                   subject,    subj_type,
                   index,      "uint:4",
                   groups,     grptype, 
                   pileup,     "log",
                   discard,    "log"
               );
               file(tf, open("fpga_types.h", "w"));
               writeln(tf, ctypes(src_type, "src_type", "SRC_"));
               writeln(tf, ctypes(subj_type, "subj_type"));
               writeln(tf, ctypes(command, "command"));
               close(tf);
            \end{lstlisting}
            
            would create file fpga\_types.h and write the following to it -
            \begin{lstlisting}[gobble=12, language=C]
                   enum src_type {
                       SRC_TIME = 0,
                       SRC_PIXEL = 1,
                       SRC_CPU = 2
                   };

                   enum subj_type {
                       NONE = 0,
                       LINEARISE = 1,
                       GAINTRIM = 2,
                       GROUP = 3,
                       PILEUP = 4,
                       THROTTLE = 5,
                       ENERGY = 6,
                       DACOLMAP = 7,
                       TIMER = 8,
                       TIMER_IM = 9
                   };

                   typedef union {
                       struct {
                           unsigned   packing: 4;
                           unsigned   discard: 1;
                           unsigned   pileup: 1;
                           unsigned   groups: 16;
                           unsigned   index: 4;
                           unsigned   subject: 4;
                           unsigned   src: 2;
                       } bits;
                       uint32_t word;
                   } command;
            \end{lstlisting}

            To allow structs packed into 32 bit words to be both transmitted
            and received as words and also accessed via struct fields it is
            best to declare a suitable union, e.g.
            \begin{lstlisting}[gobble=12, language=C]
                   typedef union {
                       struct command  assign;
                       uint32_t        word;
                   } COMMAND;
                   
                   COMMAND     c;
                   
                   c.assign.subject = ENERGY;  /* assign using fields */
                       .
                       .
                   c.assign.src = CPU;
                   
                   transmit_to_fpga(c.word);   /* transmit as 32 bit word */
            \end{lstlisting}








        \subsubsection{cdatas - generate a C name array for 3PL enum type}\label{sec:lfcdatas}

            {\bf Library: ctypes.3pl}

            {\bf Prototype:} \\
            \vsp
            \begin{lstlisting}[gobble=12, language=threepl]
               global function
               cdatas (type(t), str(name)) {}
            \end{lstlisting}

            {\bf Return mode and type:} \\
            immediate, "str"

            {\bf Input parameters:} \\            
            {\bf t} is the enum type. \\
            {\bf name} is the identifier to be used for the C string array.

            {\bf description:} \\
            This function takes a 3PL variable of type {\bf type} whose value is
            type {\bf enum} and returns a string defining a C data object that
            matches the given type. The object has the C name given by {\bf
            name}. Presently this function only operates on enum types, where
            the C object is a NULL-terminated string table formed from the enum
            values. The string array can be indexed by an enum ordinal
            to get the associated identifier string.
            
            The following 3PL code -
            \begin{lstlisting}[gobble=12, language=threepl]
               enum(e, V1, V2, V3);
               println(cdatas(e, "vname"));
            \end{lstlisting}
            would print -

            \begin{lstlisting}[gobble=12, language=C]
               static const char *vname[] __attribute((unused)) = {
                   "V1",
                   "v2",
                   "V3",
                   NULL
               };
            \end{lstlisting}



    \subsection{library classes}
    
        None.
